% Copyright (c) 2026 Geovana Neves. Licensed under CC BY 4.0 (https://creativecommons.org/licenses/by/4.0/). See LICENSE-AND-AUTHORSHIP.md.
%
% gui/text/01-reading.tex

\begin{frame}{How to read this deck}
\begin{columns}[T]
\begin{column}{0.56\textwidth}
\begin{itemize}\footnotesize
  \item A GUI session is a sequence of steps: load the geometry, mark where its
        wakes leave it, place frames and set things turning, state the flight
        condition and the solver, run, export. The deck takes them in that
        order, one table per stage.
  \item \textbf{In pyfs} says which file the key lives in: a \emph{row key}
        (a matrix column or the \code{VAR\_NAMES\_VALUES} cell), a
        \emph{setup} key, a \emph{pproc} table, the \emph{reference}, or the
        \emph{sidecar} \code{<stem>.boundaries.toml} beside the geometry.
  \item \textbf{Verified on}: the builds on which the command database
        records a probe run of every command the key writes. \emph{none}: no
        probe run, the command is still documented by its build's manual.
        \emph{no command}: the step writes none.
  \item \textbf{not yet}: no key of its own. The raw route still takes it
        (last part).
\end{itemize}
\end{column}
\begin{column}{0.40\textwidth}
{\ftstablesize\setlength{\tabcolsep}{4pt}
\begin{tabular}{@{}lrr@{}}
\toprule
\textbf{Stage} & \textbf{Key} & \textbf{Not yet} \\
\midrule
Geometry and mesh & 7 & 6 \\
Boundary conditions & 14 & 0 \\
Frames, motion, actuators & 11 & 5 \\
Flight conditions, solver & 24 & 0 \\
The run & 12 & 0 \\
Basic post & 25 & 5 \\
\midrule
\textbf{109 steps} & \textbf{93} & \textbf{16} \\
\bottomrule
\end{tabular}}
\end{column}
\end{columns}
\guisource{How to read a line}
\end{frame}
